\section{Chapter 7}

\begin{verse}
The Universe is eternal\\
It is eternal\\
Because it does not live for itself\\
Therefore: The True Man\\
By retreating, advances\\
By remaining outside, he is ever present\\
By not acting for himself, the center\\
Reaches perfection.
\end{verse}


Although the characters used in the first line are ``Heaven and Earth'', here they express the idea of totality, and it is not the ``nameable'', qualified Principle, that is referred to here but to action that makes it eternal: it ``is'' in the transcendent sense since it is denied.

The word ``retreating'' = not to put in front of himself. Reproducing in oneself the way of the Tao, the True Man advances, rises transcendentally, as the absolute individual, in so far as he does not make his human individuality the center (it can also be translated: not having person --- inclinations of person --- becomes person, he obtains a persona). Technical Taoist expressions: ``to drop the persona like a habit''.

\paragraph{Chinese text and literal translation}


\textbf{Chapter 7 (\begin{CJK*}{UTF8}{bsmi}第七章\end{CJK*})}

\columnratio{.4}
\begin{paracol}{2}
\begin{verse}
\begin{CJK*}{UTF8}{bsmi}
天長地久。\\
天地所以能長且久者，\\
以其不自生，\\
故能長生。\\
是以聖人後其身\\
而身先，\\
外其身\\
而身存。\\
非以其無私邪？\\
故能成其私。
\end{CJK*}
\end{verse}
\switchcolumn
\begin{verse}
The sky and the earth / The universe is eternal.\\
The reason for the eternal of the sky and the earth/the universe\\
Is they are not being/existing alone,\\
So they are eternal.\\
Because of this, The sages put their selves behind all other people,\\
Yet it is before all others he shall eventually stand.\\
They were indifferent to the discomforts and dangers.\\
Yet it is he who shall thus survive.\\
It is none other than that they do not demand self.\\
Yet it is he eventually who shall achieve for self.
\end{verse}
\end{paracol}

\flrightit{Posted on 2020-09-08 by Lao Tzu}

\begin{center}* * *\end{center}

\begin{footnotesize}
\begin{sffamily}
\texttt{roballenvalecourt on 2020-09-08 at 10:56 said:}

enjoying your Taoist excursions... good to see you post the Chinese characters with your interpretations... 

\hfill

\texttt{Cologero on 2020-09-08 at 19:52 said:}

Actually, the interpretations belong to Julius Evola.


\hfill

\end{sffamily}
\end{footnotesize}

